\section{Fundamente Teoretice}
\label{sec:background}

\subsection{Protocoale de tip Gossip}
Protocoalele de tip Gossip, modelează diseminarea informației într-un sistem distribuit similar modului în care se răspândește un zvon într-o rețea socială. Caracteristica definitorie a acestor protocoale este absența unei autorități centrale care să coordoneze comunicarea; în schimb, informația se propagă prin schimburi periodice între perechi de noduri.

Conform taxonomiei propuse de Shah \cite{shah2009gossip}, eficiența unui algoritm Gossip într-un mediu distribuit este condiționată de respectarea a patru principii fundamentale:
\begin{itemize} 
    \item \textbf{Localitate Strictă:} Un nod deține informații exclusiv despre vecinii săi direcți din graful de rețea. Nu este necesară menținerea unei tabele globale de rutare sau cunoașterea topologiei complete a rețelei.
    \item \textbf{Simplitate Computațională:} Algoritmul trebuie să implice operații aritmetice elementare și un consum minim de memorie, permițând rularea pe dispozitive cu resurse limitate. Interacțiunea nu trebuie să implice transferuri masive de date.
    \item \textbf{Asincronicitate:} Sistemul nu presupune existența unui ceas global sau a unor runde de comunicare perfect sincronizate. Nodurile inițiază schimburi de date independent, bazându-se pe ceasuri locale.
    \item \textbf{Toleranță la Erori (Fault Tolerance):} Algoritmul trebuie să fie robust la schimbări dinamice ale topologiei, cum ar fi defectarea nodurilor sau pierderea pachetelor, asigurând convergența rezultatului final fără intervenție externă.
\end{itemize}

\subsection{Consensul Distribuit (Average Consensus)}
Problema fundamentală abordată în această lucrare este \textit{Average Consensus}: convergența stării tuturor nodurilor către media aritmetică a valorilor inițiale din rețea. Fie $x_i(0)$ valoarea inițială a nodului $i$ (de exemplu, încărcarea CPU). Obiectivul este ca, printr-un proces iterativ, starea fiecărui nod $x_i(t)$ să satisfacă condiția \cite{spanos2005dynamic}:
\[ \lim_{t \to \infty} x_i(t) = \frac{1}{n} \sum_{j=1}^{n} x_j(0) \]

\subsection{Uniform Gossip și Conservarea Masei}
Modelul utilizat, \textbf{Uniform Gossip} \cite{kempe2003gossip}, presupune că fiecare nod selectează un partener de comunicare cu o probabilitate uniformă $1/n$ (în cazul unei topologii complet conectate). Mecanismul de actualizare se bazează pe proprietatea de \textbf{Conservare a Masei}.

Atunci când două noduri, $i$ și $j$, interacționează la momentul $t$, ele își actualizează stările la media valorilor lor curente:
\[ v_i(t+1) = v_j(t+1) = \frac{v_i(t) + v_j(t)}{2} \]
Această operație asigură invarianța sumei globale:
\[ v_i(t+1) + v_j(t+1) = v_i(t) + v_j(t) \]
Astfel, informația nu este pierdută sau creată, ci doar redistribuită, conducând sistemul către o stare de echilibru (entropie maximă).

\subsection{Dynamic Consensus}
În scenarii reale, valorile monitorizate nu sunt statice. Deoarece reinițializarea algoritmului la fiecare fluctuație a senzorilor este impracticabilă, se utilizează \textit{Dynamic Average Consensus}. Nodurile monitorizează continuu valoarea locală și, în loc să suprascrie starea algoritmului, "injectează" diferența ($\Delta$) dintre noua măsurătoare și cea anterioară direct în procesul de mediere. Aceasta permite urmărirea în timp real a mediei globale variabile \cite{spanos2005dynamic}.

\subsection{Impactul Topologiei de Rețea}
Performanța algoritmilor de consens este dictată de spectrul grafului de conexiuni:
\begin{itemize} 
    \item \textbf{Topologia Inel:} Fiecare nod are gradul $k=2$. Deși simplă, prezintă o vulnerabilitate majoră la partiționare și o viteză de convergență redusă $\bigO{N^2}$.
    \item \textbf{Topologia Stea:} Centralizează comunicarea într-un Hub. Deși diametrul este mic, Hub-ul reintroduce problema SPOF și a congestiei, contrazicând principiile distribuite.
    \item \textbf{Topologia Completă:} Permite comunicarea directă între oricare două noduri. Aceasta minimizează diametrul rețelei și maximizează viteza de amestecare, pentru că are convergență rapidă ($\bigO{\log N}$).
\end{itemize}